% SOURCE docs/release/2026-10-29/source/docs/writing-repair/2026-10-29/src/23868f4d3823c1bb.tex CHAR 0
\chapter{From a causal diagram to instrumental-variable estimation}
\label{r3:instruments}
\section{Why an invitation is not an adjustment covariate}
Suppress $X,M,C$ temporarily and consider
$Z\to A\to Y$ with $A\leftarrow U\to Y$. In a randomized invitation
study, $Z$ changes receipt while $U$ confounds receipt and outcome. The
graph expresses that $Z$ is independent of the common causes and that
all directed paths from $Z$ to $Y$ pass through $A$. Neither condition
is equivalent to $Z\perp Y\mid A$: conditioning on receipt opens the
collider path through $U$. An empirical association between $Z$ and $Y$
among activators is therefore not by itself evidence against exclusion.

The graph alone does not generally identify $\E[Y(1)-Y(0)]$
\citep[Section 3.5]{pearl2009}. Linear homogeneous effects, monotone
binary compliance, and a nonparametric structural response each add
different restrictions. This chapter develops the first two and connects
them to the inverse problem that arises in the third. Calling all three
outputs an ``IV effect'' would lose the scientific distinction.

\section{Structural exclusion becomes an observable moment}
\label{r3:linear-iv}
First suppose
\begin{equation}
 Y=\beta A+h(X)+\varepsilon,
 \qquad \E[\varepsilon\mid Z,X]=0.
 \label{r3:iv-structural}
\end{equation}
The coefficient is a homogeneous structural effect: holding the
disturbance fixed, changing $A$ by $\Delta$ changes $Y$ by $\beta\Delta$.
The structural function has no direct $Z$ argument. This exclusion and
the conditional disturbance restriction are assumptions. A structural
common cause of $A$ and $Y$ allows $\E[\varepsilon\mid A,X]\ne0$.
The linear example in \citet[Section 2]{angrist1996} starts from precisely
this distinction between endogenous receipt and exogenous assignment.

Let $r(X)=\E[A\mid X]$, $\ell(X)=\E[Y\mid X]$ and
$m(X)=\E[Z\mid X]$. Iterating the conditional disturbance restriction
gives $\ell(X)=\beta r(X)+h(X)$. Thus
\begin{align}
 0&=\E[(Z-m(X))\{Y-\ell(X)-\beta(A-r(X))\}],\nonumber\\
 \beta&=\frac{\E[(Z-m(X))(Y-\ell(X))]}
                   {\E[(Z-m(X))(A-r(X))]},
 \label{r3:residual-iv}
\end{align}
provided the denominator $\kappa$ is nonzero. Exogeneity justifies the
zero moment; relevance makes it invertible. Conditional centering removes
$h(X)$ without asserting that it is linear. A regression that uses a
finite linear basis for $X$ instead imposes the corresponding approximation
or correct-specification condition; it does not automatically perform
this population residualization.

Without covariates and with binary $Z$, the ratio reduces to
\begin{equation}
 \beta=\frac{\E[Y\mid Z=1]-\E[Y\mid Z=0]}
                   {\E[A\mid Z=1]-\E[A\mid Z=0]}.
 \label{r3:wald}
\end{equation}
Indeed, $\Cov(Z,Y)=P(Z=1)P(Z=0)\{\E[Y\mid Z=1]-\E[Y\mid Z=0]\}$,
and the same factor cancels in the receipt covariance. Nothing in this
algebra proves a homogeneous-effect model; the next section gives a
different interpretation under heterogeneous potential outcomes.

\subsection{Two-stage least squares and GMM}
For linear controls let the design matrix $D$ contain the endogenous
regressors and included controls, and let $W$ contain the instruments and
those same controls. With $W'W$ nonsingular, define
$P_W=W(W'W)^{-1}W'$. Fitting $D$ on $W$ gives $\widehat D=P_WD$;
regressing $Y$ on $\widehat D$ gives
\begin{equation}
 \widehat\beta_{\rm 2SLS}=(D'P_WD)^{-1}D'P_WY.
 \label{r3:2sls}
\end{equation}
This follows because $P_W$ is symmetric and idempotent. The structural
error used for uncertainty is $\widehat\varepsilon=Y-D\widehat\beta$,
not the residual from regressing on $\widehat D$. Treating the predicted
regressors as ordinary fixed observations with the latter residual yields
the wrong variance calculation.

More generally, define $\bar g_n(\beta)=n^{-1}\sum_iW_i(Y_i-D_i'\beta)$
and minimize $\bar g_n'\Omega_n\bar g_n$ for positive-definite
$\Omega_n$. The population derivative is $-Q_{WD}=-\E[W_iD_i']$.
With $Q_{WD}$ of full column rank, suitable moments and an independent
sample CLT, the asymptotic covariance of $\sqrt n(\widehat\beta-\beta)$ is
\begin{equation}
 (Q_{DW}\Omega Q_{WD})^{-1}Q_{DW}\Omega S\Omega Q_{WD}
 (Q_{DW}\Omega Q_{WD})^{-1},
 \label{r3:gmm-variance}
\end{equation}
where $S=\Var(W_i\varepsilon_i)$. For 2SLS,
$\Omega=\E[W_iW_i']^{-1}$; efficient GMM within this moment class uses
$\Omega=S^{-1}$. The expression follows by linearizing the moment
equation and applying the CLT to $n^{-1/2}\sum W_i\varepsilon_i$.
For clusters or time series, $S$ must be the appropriate cluster or
long-run score covariance, with the normalization and asymptotic sampling
unit adjusted accordingly. Transaction count is not the sample size in
an experiment randomized across a few branches.

\section{The same ratio can identify a complier effect}
\label{r3:late}
Now let $Z,A\in\{0,1\}$, with receipt potentials $A(0),A(1)$ and
outcome potentials $Y(z,a)$. Assume well-defined interventions, no
interference, and randomized assignment:
\[
 Z\perp\{A(0),A(1),Y(z,a):z,a\in\{0,1\}\}.
\]
Assume both assignment arms have positive probability, exclusion
$Y(z,a)=Y(a)$, and consistency
$A=A(Z)$, $Y=Y(A(Z))$. Define $\tau=Y(1)-Y(0)$.
Binary receipt gives the exact individual identity
\begin{equation}
 Y(A(1))-Y(A(0))=\tau\{A(1)-A(0)\}.
 \label{r3:compliance-identity}
\end{equation}
Taking expectations and using random assignment identifies its left side
as the outcome intention-to-treat contrast. Similarly the receipt
intention-to-treat contrast is $\E[A(1)-A(0)]$.

Let $\mathcal C=\{A(1)=1,A(0)=0\}$ be compliers and
$\mathcal D=\{A(1)=0,A(0)=1\}$ defiers. Always-takers and never-takers
have zero receipt difference. Therefore the Wald functional is
\begin{equation}
 \frac{P(\mathcal C)\E[\tau\mid\mathcal C]
             -P(\mathcal D)\E[\tau\mid\mathcal D]}
       {P(\mathcal C)-P(\mathcal D)}.
 \label{r3:iv-with-defiers}
\end{equation}
Monotonicity $A(1)\ge A(0)$ removes defiers. Relevance then implies
$P(\mathcal C)>0$, yielding the local average treatment effect
\begin{equation}
 \theta_{\mathcal C}=\E[Y(1)-Y(0)\mid\mathcal C]
 =\frac{\Delta_Z Y}{\Delta_Z A}.
 \label{r3:late-result}
\end{equation}
This reconstructs \citet[Equation (9), Table 1 and Proposition 1]{angrist1996}.
Monotonicity is not an arrow property. An invitation can discourage some
customers, for example by attracting attention to a competing product.
Without monotonicity, Equation~\eqref{r3:iv-with-defiers} can lie outside
the range of subgroup average effects.

Compliance strata refer to joint receipt potentials, not observed groups.
An observed nonactivator under $Z=0$ may be a complier or never-taker.
The identified average effect therefore does not label individual
compliers. Changing the invitation, price or delivery channel can change
the complier population and its mean effect. LATE is not automatically
the gain from expanding the product to all eligible households.

\subsection{A constructive failure to identify the population ATE}
Take $40\%$ compliers, $30\%$ always-takers and $30\%$ never-takers,
with fair randomized $Z$. Among compliers set $(Y(0),Y(1))=(0,1)$.
Always-takers have observed $Y(1)=1$; never-takers have observed $Y(0)=0$.
Consider two completions of their missing potential outcomes:
\begin{center}
\begin{tabular}{lcc}
\toprule
Stratum & Completion I $(Y(0),Y(1))$ & Completion II $(Y(0),Y(1))$\\
\midrule
Complier & $(0,1)$ & $(0,1)$\\
Always-taker & $(1,1)$ & $(0,1)$\\
Never-taker & $(0,0)$ & $(0,1)$\\
\bottomrule
\end{tabular}
\end{center}
Both satisfy exclusion and monotonicity and yield the same full observed
law of $(Z,A,Y)$: receipt rates are $0.3$ and $0.7$, with $Y=A$.
The first stage and reduced form are both $0.4$, so LATE is one. But
the ATE is $0.4$ in completion I and one in completion II. The observed
IV experiment cannot distinguish them. This is a project counterexample,
not an appeal to a generic statement that ``IV has limitations.'' It
identifies the exact missing outcomes that a population effect requires.

\subsection{Conditional instruments and the weights of the target}
If assignment is randomized only conditional on $X$, assume the IV
conditions within each $x$, and let $p(x)=P(Z=1\mid X=x)$,
$\Delta_A(x)=\E[A\mid Z=1,x]-\E[A\mid Z=0,x]$ and
$\tau_{\mathcal C}(x)$ be the conditional complier effect. Then
\[
 \frac{\E[(Z-p(X))Y]}{\E[(Z-p(X))A]}
 =\frac{\E[p(X)(1-p(X))\Delta_A(X)\tau_{\mathcal C}(X)]}
        {\E[p(X)(1-p(X))\Delta_A(X)]}.
\]
The identity follows by computing each conditional covariance. Even with
monotonicity, the residual-IV ratio weights strata by both assignment
variance and compliance. A policy that targets a different $X$ distribution
or uses different invitation probabilities generally changes the target.
The product must preserve these weights rather than translate the result
into an unqualified population ATE.

Transporting a complier effect needs more than target covariate weights.
Let $\kappa^*(x)=P^*\{A(1)>A(0)\mid X=x\}$ be the target-domain
compliance probability. Its overall complier effect is
\[
 \tau_{\mathcal C}^*
 =\frac{\E^*[\kappa^*(X)\tau_{\mathcal C}^*(X)]}
        {\E^*[\kappa^*(X)]}.
\]
Replacing the conditional complier effect by its source estimate
requires invariance of that complier response, a stronger restriction
than sharing an outcome name. Target compliance must also be identified,
for example from a valid target encouragement design, and have source
support. Different pricing or eligibility rules can change who complies.
An overall source Wald ratio and target demographics alone do not
identify this expression. This is a target-specific consequence of the
compliance decomposition, not a generic extrapolation guarantee.

\section{Weak and invalid instruments fail differently}
\label{r3:weak-iv}
In the scalar linear model write $N=\E[(Z-m)(Y-\ell)]$ and
$\kappa=\E[(Z-m)(A-r)]$, with $N=\beta\kappa$.
For fixed nonzero $\kappa$, a ratio expansion gives
\[
 \widehat\beta-\beta
 =\frac{(\widehat N-N)-\beta(\widehat\kappa-\kappa)}{\kappa}
  +o_p(n^{-1/2}).
\]
If $\kappa$ shrinks on the same scale as its sampling error, the
denominator is no longer stable and this expansion fails. Orthogonal
nuisance learning can reduce error in $\widehat N$ and
$\widehat\kappa$; it cannot create population relevance.

The local-to-zero experiment makes the failure precise. Suppress controls,
normalize $\E[Z^2]=1$, and let
$A=\pi_n Z+v$, $Y=\beta A+u$, $\pi_n=c/\sqrt n$.
Suppose the sample second moment converges to one and
$n^{-1/2}\sum_i(Z_i u_i,Z_i v_i)\Rightarrow(G_u,G_v)$, a
nondegenerate joint Gaussian vector. Then
\begin{equation}
 \widehat\beta-\beta
 =\frac{n^{-1/2}\sum_i Z_i u_i}
 {c\,n^{-1}\sum_iZ_i^2+n^{-1/2}\sum_iZ_i v_i}
 \ \Longrightarrow\ \frac{G_u}{c+G_v}.
 \label{r3:weak-limit}
\end{equation}
The error does not concentrate at zero, even as $n$ grows. This is the
scalar specialization of the local-to-zero construction in
\citet[Section 2A, equations (2.4)--(2.7)]{staiger1994}. The number of
customers alone therefore cannot certify that a weak encouragement
supports precise receipt-effect inference.

\subsection{Test the structural residual before dividing}
Let $X$ now denote a full-rank matrix of $k$ included controls, $Z$ a
matrix of $q$ excluded instruments, and $M_X=I-P_X$. Suppose
$[X,Z]$ has rank $k+q$. Set $\widetilde Z=M_XZ$,
$P_{\widetilde Z}=\widetilde Z(\widetilde Z'\widetilde Z)^{-1}
\widetilde Z'$ and $r_b=M_X(Y-bA)$. The Anderson--Rubin construction is
\begin{equation}
 F_{\rm AR}(b)=
 \frac{r_b'P_{\widetilde Z}r_b/q}
      {r_b'(M_X-P_{\widetilde Z})r_b/(n-k-q)}.
 \label{r3:ar-statistic}
\end{equation}
Under $b=\beta$, $r_b=M_Xu$. If conditional on $X,Z$ the structural
error is $N(0,\sigma^2I)$, the two quadratic forms divided by
$\sigma^2$ are independent chi-squared variables of ranks $q$ and
$n-k-q$: their projection spaces are orthogonal. Thus
$F_{\rm AR}(\beta)\sim F_{q,n-k-q}$ without a lower bound on the
first stage. Invert this test to form
$\{b:F_{\rm AR}(b)\le F_{q,n-k-q;1-\alpha}\}$.
The derivation reconstructs the exact Gaussian result reported in
\citet[Section 2D, equation (2.24) and Theorem 3]{staiger1994}; that
paper separately supplies an asymptotic result under its moment
conditions. Gaussian, homoskedastic fixed-design inference is not a
cluster-robust theorem.

With heteroskedastic or dependent observations, a candidate-specific
moment statistic can instead use
$n\bar m(b)'\widehat S(b)^{-1}\bar m(b)$, where
$m_i(b)=\widetilde Z_i(Y_i-bA_i-X_i'\widehat\gamma_b)$.
Its reference distribution requires a valid score CLT, a consistent
covariance and fixed-rank restrictions; estimated controls and learned
instruments need their own remainder analysis. Uniform validity across
weak first stages must be established for that construction, not inferred
from the Gaussian calculation. Independent training of a compressed
instrument can avoid using the same errors to select and test it;
the split-sample AR discussion in \citet[printed p.17]{staiger1994}
provides an early instance, with independence assumptions explicit.

\subsection{Why the answer can be disconnected}
For an elementary exact Gaussian illustration, suppose a reduced-form
statistic $(\widehat N,\widehat\kappa)$ has mean
$(\beta\kappa,\kappa)$ and known covariance $V$. Testing
$\widehat N-b\widehat\kappa$ gives the confidence-set inequality
\begin{equation}
 (\widehat\kappa^2-z^2V_{\kappa\kappa})b^2
 +(-2\widehat N\widehat\kappa+2z^2V_{N\kappa})b
 +(\widehat N^2-z^2V_{NN})\le0,
 \label{r3:ratio-set}
\end{equation}
where $z=z_{1-\alpha/2}$. This follows by squaring a normal test of a
linear combination; it is a project calculation in a specified Gaussian
experiment, not a claim that estimated sample covariances are exact.
When the leading coefficient is negative and the discriminant positive,
the accepted set consists of two exterior rays. It can also be all of
$\mathbb R$. Replacing that result by the smallest finite interval on
a plotting grid would discard the scientific content: the experiment
does not rule out very large effects. A bank may then prefer the
well-identified invitation ITT or a stronger encouragement design.
The same inverse structure reappears when uncertain entity linkage mixes two
portfolio responses. Section~\ref{s61:linkage-sets} diagonalizes that map and
constructs financial confidence sets from candidate responses and audit
counts, retaining rank failure and disconnected answers.

Invalidity is a different problem. For centered scalar variables,
suppose $A=\pi Z+v$ and $Y=\beta A+\delta Z+u$, with $Z$ orthogonal
to both $u,v$. Then
\begin{equation}
 \frac{\Cov(Z,Y)}{\Cov(Z,A)}=\beta+\frac{\delta}{\pi}.
 \label{r3:invalid-iv}
\end{equation}
The instrument may be strongly predictive and precisely estimated while
the exclusion is false. A direct invitation effect of $0.02$ outcome units
produces bias $0.05$ when $\pi=0.4$, but bias $1$ when $\pi=0.02$.
Weakness amplifies a small violation. An overidentification test can reject
a joint set of restrictions; failure to reject does not prove that all
instruments satisfy exclusion, especially when violations are shared.

For a declared sensitivity range $|\delta|\le d$ in this scalar model,
the admissible coefficient interval is centered at the covariance ratio
with radius $d/|\pi|$, before sampling uncertainty is added. That interval
comes from a substantive bound on a direct effect, not from a diagnostic
threshold chosen by software.

\section{Machine learning enters after the moment restriction}
\label{r3:iv-ml}
For Equation~\eqref{r3:iv-structural}, use the score
\begin{equation}
 \psi(O;\beta,\eta)=(Z-m(X))
       \{Y-\ell(X)-\beta(A-r(X))\},\quad \eta=(m,\ell,r).
 \label{r3:pliv-score}
\end{equation}
At the truth its derivatives in arbitrary square-integrable nuisance
directions vanish: an $\ell$ or $r$ perturbation multiplies a function
of $X$ by $Z-m(X)$, whose conditional mean is zero; an $m$ perturbation
multiplies a function of $X$ by $\varepsilon$, whose conditional mean
is zero. The derivative in $\beta$ is $-\kappa$, which must still be
bounded away from zero for conventional strong-identification inference.
This is the partially linear IV score of
\citet[Section 4.2, Equations (4.5)--(4.8)]{dml2018}.

Estimate $m,\ell,r$ on other folds, evaluate residual products on the
held-out units, and solve the sample moment. The nuisance contribution
to its population bias is
\[
 \E[(\widehat m-m)\{(\widehat\ell-\ell)
                         -\beta(\widehat r-r)\}],
\]
because the first-order terms have zero expectation. Cauchy--Schwarz
bounds this by a product of nuisance errors. The useful ML rates are
therefore rates for these moment errors and the score variance, not
merely predictive accuracy for $Y$. Dependence determines valid fold
construction and the CLT; arbitrarily splitting rows from the same
customer does not produce independent training and evaluation samples.

\section{Control functions use a triangular structural restriction}
\label{r3:control-function}
There is another IV route that must not be confused with adjusting
for an arbitrary first-stage residual. Suppose
\[
 A=h(Z,X,V),\qquad Y=f(A,X,U),\qquad
 Z\perp(U,V)\mid X,
\]
where $V$ is scalar, continuously distributed and $h(z,x,v)$ is
strictly increasing in $v$. Normalize $V\mid X$ to be uniform on
$(0,1)$. Then the observed conditional rank
\begin{equation}
 R=F_{A\mid Z,X}(A\mid Z,X)=V
 \label{r3:control-rank}
\end{equation}
almost surely. For fixed $z,x$, the event
$h(z,x,V)\le h(z,x,v)$ is $V\le v$; independence of $Z$ proves
the equality. Conditional on $X,R$, the dependence between $U$ and
the endogenous disturbance is held fixed and the remaining action
variation comes from $Z$, so $U\perp A\mid X,R$. Consequently
\[
 m_{\rm cf}(a,x,r)=\E[Y\mid A=a,X=x,R=r]
                 =\E[f(a,x,U)\mid X=x,V=r]
\]
is identified on the joint observed support. This reconstructs
\citet[Assumptions 3.1--3.2, Theorem 1 and Appendix A]{imbens2002},
with observed covariates added by applying the argument within their
strata. Strict monotonicity in the latent rank is distinct from
LATE's monotonic response to the observed instrument.

To recover the average structural response at a specified $a$, one
would integrate
\begin{equation}
 \E[Y(a)]=\int\!\int_0^1m_{\rm cf}(a,x,r)\,dr\,dP_X(x).
 \label{r3:control-average}
\end{equation}
This requires $a$ to have support at almost every rank and covariate
being integrated. The source's Assumption 3.3 and Theorem 2 make this
additional support condition explicit. First-stage predictability
or rank recovery alone does not supply it.

For an exact counterexample, let $Z,V$ be independent uniform
variables on $[0,1]$ and $A=Z+V$. The two response systems
\[
 f_1(a,v)=a+v,\qquad
 f_2(a,v)=a+v+c\ind\{a<v\}
\]
agree on every observed record because $A\ge V$. Both obey the
triangular model and strict first-stage monotonicity. Under
$\doop(A=.2)$, however, their average responses differ by $.8c$.
At $A=.2$ the observed ranks cannot exceed $.2$, while the intervention
preserves all ranks. This is failure of structural-response
identification caused by support, despite knowing the control exactly.

In the additive first stage $A=\pi Z+V$, with
$\E[U\mid V]=\lambda V$ and $Y=\theta A+U$, regressing on
$A$ and the true residual $V$ is a sufficient special case. Outside
it, a linear residual adjustment can miss a nonlinear control
function or fail because several latent dimensions determine action.
For discrete $A$, the conditional CDF is generally not an invertible
latent rank. Estimated ranks are generated regressors: inference must
account for the first-stage CDF estimation and any support trimming.
The source's two-step series estimator is one qualified route; it
does not establish validity for an arbitrary learned encoder called
a control function.

\section{A nonlinear structural response is an inverse problem}
\label{r3:inverse-introduction}
Suppose the model is $Y=h(A,X)+\varepsilon$ with
$\E[\varepsilon\mid Z,X]=0$. At fixed $x$ the observable regression
$g(z,x)=\E[Y\mid z,x]$ satisfies
\begin{equation}
 (T_xh_x)(z):=\int h(a,x)\,dP(a\mid z,x)=g(z,x).
 \label{r3:npiv-operator}
\end{equation}
The operator $T_x$ averages a structural response over the first-stage
conditional distribution. Identifying $h_x$ requires injectivity on its
declared function class: $T_xv=0$ must imply $v=0$. In unrestricted
function classes this is a completeness condition, far stronger than a
nonzero scalar first-stage covariance. A binary instrument supplies only
two conditional expectations at a given $x$; those alone cannot identify
an arbitrary function of continuous $A$.

Even an injective operator may be ill conditioned. For a singular system
$T\phi_j=\sigma_j\chi_j$ with small $\sigma_j$, inversion divides noisy
coefficients by $\sigma_j$. Ridge regularization instead estimates
\[
 h_\lambda=\sum_j\frac{\sigma_j}{\sigma_j^2+\lambda}
              \langle g,\chi_j\rangle\phi_j.
\]
It suppresses unstable directions and introduces bias. The formula is a
direct singular-value calculation; a full statistical NPIV survey must
specify source conditions, estimation of $T$, rates and uncertainty for
the actual target functional. Deep networks are function approximators
inside this inverse problem; their expressiveness does not establish
injectivity or valid exclusion.
\citet[Sections 2 and 4]{hall2005} gives kernel and series estimators
with rate analysis for an estimated operator. Chapter~\ref{r3:inverse-learning}
develops the distinction between function and functional identification,
regularization and proxy bridges, rather than treating inverse learning
as ordinary response regression.

A control-function route uses different assumptions. If
$A=f(Z,X,V)$ is strictly increasing in scalar $V$, has continuous
conditional distributions and $(V,\varepsilon)\perp Z\mid X$, then
$F_{A\mid Z,X}(A\mid Z,X)=F_{V\mid X}(V\mid X)$ is a control variable.
Conditioning on it and $X$ can remove dependence induced by $V$, subject
to support and a suitable outcome structure. This rank construction
follows by applying the monotone inverse of $f$. Binary receipt does not
have an invertible rank of this form; inserting a binary first-stage
residual into an arbitrary outcome model is not justified by the argument.

Proxy approaches solve related conditional expectation equations but
assign different roles to their variables. A variable measuring latent
liquidity need can be a useful proxy while being an invalid instrument
because it predicts outcome through that need. Instrumental, proxy,
adjustment and learned-representation roles therefore require separate
measurement and exclusion assumptions. The theoretical survey must
compare them at those restrictions, not at their model architectures.

\section{Additional assumptions change the answer}
\label{r3:iv-contract}
The invitation experiment can return an intention-to-treat effect under
randomization. A homogeneous linear model can return a structural
coefficient and moment-based uncertainty under rank conditions. Monotone
binary compliance can return a complier effect and its population weights.
An unrestricted IV graph can instead return bounds or an explicit
nonidentification result. A regularized inverse problem returns an estimate
conditional on a function class, operator restrictions and regularization.

Each result must carry the original intervention, population, assumptions,
first-stage evidence, score or likelihood, dependence unit and uncertainty
method. A weak first stage should affect the inference state; a disputed
exclusion should affect the causal assumption state. Combining them into
one opaque confidence score would conceal which remedy is required.
Additional randomization can improve support; a different encouragement
may improve relevance; measuring an omitted mechanism may alter the graph.
None of these changes is equivalent to running a more accurate predictor.

